\clearpage

\begin{center}
    \section*{Аннотация}
    % \addcontentsline{toc}{section}{Аннотация}
\end{center}

\begin{center}
    \section*{Abstract}
    % \addcontentsline{toc}{section}{Abstract}
\end{center}

\begin{center}
    \section*{Введение}
    \addcontentsline{toc}{section}{Введение}
\end{center}


В эпоху информационной перегрузки
возможность эффективно управлять данными
стала особенно важной.
Классических способов работы с информацией,
таких как ведение заметок,
уже недостаточно,
особенно для специалистов
регулярно работающими с огромными массивами многотемных
данных. В качестве решения были разработаны специализированные системы
для управления знаниями, например, Obsidian.
Они берут на себя роль «второго мозга»,
освобождая пользователя от хранения всей информации у себя в голове
и упрощая манипуляции с ней.

Современные системы управления знаниями поддерживают
генерацию с дополненной выборкой (Retrieval-Augmented Generation, RAG).
Большинство RAG-систем используют гибридный поиск
(комбинацию лексического и семантического).
Такой подход имеет несколько слабостей:
лексический поиск не умеет искать синонимы,
а семантический может упустить релевантные данные из-за большой разницы
в векторных представлениях (эмбеддингах).

Более продвинутой стратегией поиска является
использование графов знаний (ГЗ),
которые соединяют семантически связанные сущности.
Каждая связь имеет направление и тип,
отражающий семантическую связь.
Это открывает доступ к многошаговому поиску,
что решает проблемы гибридного поиска и
позволяет увеличить полноту результатов.

Несмотря на преимущества ГЗ,
их построение и поддержка требуют колоссальных
ресурсов. По мере роста базы знаний возникает «узкое горлышко»:
пользователям недостаточно когнитивных ресурсов и времени
на анализ и связывание заметок.

На данный момент существующие решения
(графовые базы данных корпоративного уровня или
облачные приложения )
имеют критические недостатки.
Базы данных требуют дорогостоящей инфраструктуры
и чрезмерно зависят от больших языковых моделей (LLM),
что приводит к генерации бессмысленных связей и галлюцинациям.
Облачные решения, несмотря на их удобство для повседневного использования,
создают риск утечки данных и принудительно привязывают пользователей к поставщику.

Цель работы -- разработка автоматизированного алгоритма
для построения графов знаний путем семантического связывания заметок в экосистеме Obsidian,
а также сравнение разработанного решения с существующими аналогами в сценариях использования RAG.

Для достижения поставленной цели были сформулированы следующие задачи:

\begin{enumerate}
    \item Провести анализ существующих решений для построения графов знаний
          для выявления их ограничений и функциональных особенностей.
    \item Разработать алгоритм автоматического поиска и типизации связей, включающий:
          извлечение содержимого Markdown-файлов; гибридный поиск кандидатов; классификацию связей; сохранение
          найденных отношений, поддерживающее инкрементальные обновления.
    \item Создать систему валидации и репрезентативный аннотированный набор данных для объективной оценки качества работы алгоритма
          и его сравнения с аналогами.
    \item Провести итеративное улучшение решения путем экспериментов с различными моделями,
          гиперпараметрами и методами фильтрации связей
          для повышения как точности системы, так и ее вычислительной эффективности.
\end{enumerate}

Объект исследования — процессы поиска,
структурирования и извлечения информации в системах управления знаниями.
Предмет исследования — методы и алгоритмы автоматического
семантического связывания текстовых данных для построения ГЗ.

Алгоритм, полученный в результате данной работы,
ориентирован на продвинутых пользователей баз знаний,
которым может быть важна суверенность данных,
возможность работы алгоритма локально
или при ограниченных вычислительных мощностях,
а также эффективно обновлять ГЗ.
Это подтверждает практическую значимость работы.

% МОЙ ТЕКСТ
% Введение
% (Пишется обычно в самом конце, но каркас нужен сразу)
%     Актуальность темы: Информационная перегрузка, рост PKM-систем, ограничения гибридного поиска в RAG и сложность ручного ведения Knowledge Graphs (KG).
%     Цель работы: Разработка автоматизированного пайплайна для извлечения, фильтрации и типизации связей между заметками в экосистеме Obsidian.
%     Задачи: (Прямо из твоего пропоузала: анализ решений, разработка алгоритма, создание датасета и фреймворка оценки, тестирование).
%     Объект и предмет исследования.

% Глава 1. Теоретические основы и обзор предметной области (Теория)
% Здесь мы раскрываем твои пункты 1, 2, 3 и 4. Глава должна плавно подвести читателя к тому, что существующие решения не идеальны, и твой подход необходим.
%     1.1. Управление персональными знаниями (PKM) и графы знаний
%         Концепция "Второго мозга" и экосистема Obsidian.
%         Роль Knowledge Graphs (KG) в структурировании информации.
%         Проблема "бутылочного горлышка" при ручном ведении масштабных PKG (Personal Knowledge Graphs).
%     1.2. Технологии поиска и обработки естественного языка в RAG
%         Ограничения классического лексического (BM25) и семантического (Cosine Similarity) поиска. Концепция гибридного поиска.
%         Применение Retrieval-Augmented Generation (RAG) в системах управления знаниями.
%         Обзор языковых моделей: архитектура трансформеров, LLM (напр., Gemini, Llama) и легковесные NLI-модели (Cross-encoders, DeBERTa).
%     1.3. Анализ существующих решений для построения графов знаний
%         Enterprise-решения: Neo4j, Memgraph (плюсы: масштаб, минусы: статика, оверхед, галлюцинации LLM).
%         Академические пайплайны: EduKG (плюсы: микро-графы, эффективность, минусы: привязка к онтологиям).
%         Облачные PKM и плагины: Mem.ai, Smart Connections, Khoj (плюсы: интеграция, минусы: vendor lock-in, "черный ящик", отсутствие сложной семантики).
%     1.4. Методы оценки качества графов и систем RAG
%         Метрики классификации для проверки качества связей (Precision, Recall, F1-score).
%         Фреймворки оценки LLM и RAG-систем (в частности, RAGAS: Context Precision, Answer Faithfulness).



% Глава 3. Экспериментальное исследование и оценка результатов (Практика: Часть 2)
% Здесь ты доказываешь, что твой код работает лучше аналогов.
%     3.1. Дизайн эксперимента и Baseline-решение
%         Описание тестового стенда и гиперпараметров.
%         Создание базового решения (baseline) на основе простого dense retrieval + LLM для сравнения.
%     3.2. Оценка качества извлечения и типизации связей
%         Сравнение пайплайна с baseline по метрикам Precision и Recall на собранном датасете.
%         Анализ ошибок (где алгоритм ошибся и почему).
%     3.3. Эффективность построенного графа в задачах RAG
%         Оценка влияния сгенерированных графов на качество ответов RAG.
%         Использование RAGAS (Context Precision, Answer Faithfulness).
%     3.4. Оценка вычислительной эффективности
%         Анализ затрат (скорость работы локальных моделей vs. стоимость/скорость API).
%         Ограничения разработанного решения.

% Заключение
%     Краткие выводы по каждой главе.
%     Подтверждение того, что цель достигнута (пайплайн создан, когнитивная нагрузка снижена, RAG улучшен).
%     Направления для дальнейших исследований (Future work).

% Список литературы
% (Сюда отлично лягут все твои источники: Lewis, Gao, Edge, Nogueira, He и документации к Neo4j/LanceDB).


% % Объект работы
% Объект работы --- мультимодальное машинное обучение.

% % предмет работы
% Предмет работы --- применение мультимодального машинного обучения к задаче
% классификации вида больной/здоровый человек.

% % Цель работы
% Цель работы --- изучение применимости мультимодального машинного обучения на основе структурированных
% и неструктурированных данных к задаче классификации больных и здоровых людей.

% % Задачи
% \begin{itemize}
%     \item[] Задачи:
%           \begin{itemize}
%               \item[$\bullet$] Изучить текущий прогресс машинного обучения в области медицины
%               \item[$\bullet$] Изучить процесс разработки мультимодальной модели
%               \item[$\bullet$] Изучить виды слияния модальностей
%               \item[$\bullet$] Найти мультимодальный датасет, состоящий из
%                     структурированных и неструктурированных данных
%               \item[$\bullet$] Изучить имеющиеся решения на выбранном датасете
%               \item[$\bullet$] Выбрать решение, предложить к нему улучшения,
%                     провести над ним эксперименты
%               \item[$\bullet$] Проанализировать результаты, сделать выводы
%           \end{itemize}
% \end{itemize}
